\section{Methodology}
\label{sec:method}

\subsection{System Overview}

Figure~\ref{fig:pipeline} illustrates the three-phase pipeline from development artifacts to research papers.

\begin{figure}[t]
\centering
\begin{verbatim}
  ┌──────────────────┐
  │   Development    │ Commits, waves, docs, roadmaps
  │    Artifacts     │
  └────────┬─────────┘
           ↓
  ┌──────────────────┐
  │ Phase 1: Topic   │ Cluster artifacts by theme,
  │   Discovery      │ score by novelty/evidence/fit
  └────────┬─────────┘
           ↓
  ┌──────────────────┐
  │ Phase 2: Evidence│ Extract quantitative data,
  │   Extraction     │ design rationale, references
  └────────┬─────────┘
           ↓
  ┌──────────────────┐
  │ Phase 3: Paper   │ Generate LaTeX structure,
  │   Generation     │ write 6 sections, compile
  └────────┬─────────┘
           ↓
  ┌──────────────────┐
  │ Research Paper   │ main.tex, sections/, PDF
  └──────────────────┘
\end{verbatim}
\caption{Three-phase pipeline from development artifacts to research papers.}
\label{fig:pipeline}
\end{figure}

\subsection{Phase 1: Topic Discovery}

Topic discovery scans development artifacts to identify research-worthy themes.

\textbf{Inputs}:
\begin{itemize}
\item Git commit history (git log --oneline --all)
\item Project waves (wave-index.yaml, wave metadata, pulses)
\item Design documents (DESIGN.md, RFC files, ADR records)
\item Roadmaps (ROADMAP.md with priority sections)
\end{itemize}

\textbf{Process}:
\begin{enumerate}
\item \textbf{Artifact scanning}: For each source:
  \begin{itemize}
  \item \textbf{Commits}: Parse commit messages, extract keywords, group by prefix ([panel], [waves]).
  \item \textbf{Waves}: Load wave-index.yaml, extract wave titles, descriptions, and pulses.
  \item \textbf{Docs}: Parse markdown files for headings, technical terms, design rationale.
  \item \textbf{Roadmaps}: Extract "Not started" items from priority sections with venue targets.
  \end{itemize}

\item \textbf{Clustering}: Group artifacts by theme using:
  \begin{itemize}
  \item Keyword matching (e.g., "hierarchical", "three-tier", "board" → cluster as "hierarchical review architecture")
  \item Temporal proximity (commits within same week likely related)
  \item Wave membership (pulses in same wave likely address same theme)
  \end{itemize}

\item \textbf{Scoring}: For each cluster, compute novelty, evidence, and fit scores:
  \begin{itemize}
  \item \textbf{Novelty} (0-1): First implementation of pattern (1.0), incremental improvement (0.5), bug fix (0.1)
  \item \textbf{Evidence} (0-1): Quantitative data available (1.0), qualitative observations (0.5), no data (0.0)
  \item \textbf{Fit} (0-1): Match to research venues (ICSE/CHI/MSR): architectural contribution (1.0), empirical study (0.8), tool demo (0.6)
  \end{itemize}
  Total score = 0.4 × novelty + 0.4 × evidence + 0.2 × fit

\item \textbf{Proposal generation}: For each cluster with score ≥ 0.6:
  \begin{itemize}
  \item Generate paper title (e.g., "Hierarchical Review Architecture: Scaling Expert Feedback Through Three-Tier Synthesis")
  \item Suggest venue target based on contribution type
  \item Summarize evidence (commit count, wave count, quantitative data available)
  \end{itemize}
\end{enumerate}

\textbf{Outputs}: List of proposed papers with title, venue, evidence summary, and score.

\subsection{Phase 2: Evidence Extraction}

Evidence extraction locates and structures data relevant to each approved paper topic.

\textbf{Inputs}:
\begin{itemize}
\item Approved paper topic (title, venue, theme)
\item Development artifact sources (as in Phase 1)
\end{itemize}

\textbf{Process}:
\begin{enumerate}
\item \textbf{Relevance filtering}: For each artifact, compute relevance to topic using keyword matching and semantic similarity (sentence embeddings, cosine similarity threshold 0.7).

\item \textbf{Quantitative extraction}: Scan relevant artifacts for numbers:
  \begin{itemize}
  \item Commit messages: "47/47 tests passed" → extract 47, 47, 100\%
  \item Wave descriptions: "14 papers across 3 modules" → extract 14, 3
  \item Review files: "Score: 3/4" → extract score data
  \end{itemize}

\item \textbf{Qualitative extraction}: Extract design rationale from:
  \begin{itemize}
  \item Commit messages (motivation for change)
  \item Wave descriptions (problem being solved)
  \item Design documents (architectural decisions)
  \end{itemize}

\item \textbf{Timeline construction}: Order relevant commits and waves chronologically to show evolution of the system.

\item \textbf{Cross-referencing}: Link artifacts (e.g., commit 6b9286e implements Wave 5 pulse h8i9j2 "active-revisions").
\end{enumerate}

\textbf{Outputs}:
\begin{itemize}
\item Evidence database: structured records linking artifacts to paper sections
\item Quantitative data table: numbers ready for results section
\item Timeline: chronological development narrative
\end{itemize}

\subsection{Phase 3: Paper Generation}

Paper generation creates LaTeX structure and writes substantive content drawing on extracted evidence.

\textbf{Inputs}:
\begin{itemize}
\item Paper topic (title, venue, theme)
\item Evidence database from Phase 2
\item Paper templates (structure, citation style, section outlines)
\end{itemize}

\textbf{Process}:

\textbf{3.1: Structure generation}
\begin{enumerate}
\item Create directory: research/panel-\{slug\}/
\item Generate main.tex with preamble, abstract, section includes, bibliography
\item Create sections/ directory with 6 files:
  \begin{itemize}
  \item 01-introduction.tex
  \item 02-related-work.tex
  \item 03-methodology.tex
  \item 04-results.tex
  \item 05-discussion.tex
  \item 06-conclusion.tex
  \end{itemize}
\item Create Makefile for building PDF (pdflatex + bibtex)
\item Create \_panel.yaml with initial state (stage: draft, round: 0)
\end{enumerate}

\textbf{3.2: Content generation by section}

\textbf{Introduction (01-introduction.tex)}:
\begin{itemize}
\item Problem statement: What gap does this work address?
\item Motivation: Why is the problem important? (reference evidence timeline)
\item Contributions: List 3-5 key contributions (reference quantitative findings)
\item Paper outline: Briefly describe remaining sections
\end{itemize}

\textbf{Related work (02-related-work.tex)}:
\begin{itemize}
\item Survey 4-6 related areas (mining software repositories, automated documentation, etc.)
\item For each area, cite 2-4 papers and explain relationship to this work
\item Conclude with gap analysis: what existing work doesn't address
\end{itemize}

\textbf{Methodology (03-methodology.tex)}:
\begin{itemize}
\item System architecture: Describe phases/components (reference design documents)
\item Implementation details: Algorithms, data structures (reference commits)
\item Design decisions: Justify choices (reference wave descriptions)
\end{itemize}

\textbf{Results (04-results.tex)}:
\begin{itemize}
\item Evaluation setup: Dataset, metrics, baseline
\item Quantitative findings: Tables and figures with extracted numbers
\item Qualitative findings: Observations from development history
\end{itemize}

\textbf{Discussion (05-discussion.tex)}:
\begin{itemize}
\item Interpretation of results: Why did numbers come out this way?
\item Limitations: What doesn't the system handle? (reference "requires human judgment" evidence)
\item Generalization: How does this apply to other domains?
\end{itemize}

\textbf{Conclusion (06-conclusion.tex)}:
\begin{itemize}
\item Summary: Restate contributions
\item Future work: Next steps (reference roadmap items)
\item Acknowledgments: Contributors, funding
\item AI disclosure: Note that paper was generated by AI from artifacts
\end{itemize}

\textbf{3.3: Compilation and validation}
\begin{enumerate}
\item Run make pdf to compile LaTeX
\item Check for errors (missing references, syntax errors)
\item If errors: fix automatically where possible, flag for human review otherwise
\item Generate PDF output
\end{enumerate}

\textbf{Outputs}:
\begin{itemize}
\item Complete paper directory with LaTeX source
\item Compiled PDF (if validation successful)
\item \_panel.yaml state file for tracking review progress
\end{itemize}

\subsection{Self-Referential Paper Generation}

This paper (panel-meta-research-automation) was generated using the system described above:

\begin{enumerate}
\item \textbf{Topic discovery}: Scanned 6 waves of panel development, identified "meta-research automation" theme (Wave 4 pulse g7h8ia, commit 7eac5de).

\item \textbf{Evidence extraction}: Found panel:import implementation, topic-discovery.md and paper-generator.md modules, 5 papers generated using the system.

\item \textbf{Paper generation}: Created this LaTeX structure, wrote 6 sections drawing on wave descriptions and commit messages, compiled to PDF.
\end{enumerate}

This creates a self-referential loop: the system documents its own methodology, providing both technical contribution and meta-level validation.

\subsection{Implementation}

The system is implemented as part of the panel plugin for Claude Code:

\begin{itemize}
\item \textbf{Topic discovery}: panel:import --from waves, --from commits, --from roadmap
\item \textbf{Evidence extraction}: shared/topic-discovery.md module
\item \textbf{Paper generation}: shared/paper-generator.md module
\end{itemize}

The system uses Claude (Anthropic) for natural language generation, git for commit history access, and LaTeX for document compilation.